\chapter{Das dauernde Ich}

\section{Drei Zeiten}

Baumann hat, um die Zeit \zitat{nach ihrer eigenen Art} zu behandeln, drei Schichten unterschieden. Die erste ist die psychologische Zeit. Das bloße Nacheinander von Vorstellungen wäre noch keine Zeitvorstellung. \zitat{Es gehört zur Zeitvorstellung ausser dem Nacheinander der Vorstellungen etwas, das sich dieses Nacheinanders als solchen bewusst wird [\ldots], etwas Bleibendes in der Aufeinanderfolge der Ideen. Dies Bleibende ist in uns unsere Ichvorstellung.}\footnote{Baumann, Lehren, Bd.\,2, S.\,659.} Die Zeit entsteht, indem das Nacheinander auf das dauernde Ich bezogen wird, und die Dauer des Ich ist selbst nicht Zeit; sie liegt der Zeit voraus, näher der Ewigkeit als der Zeit. Kants Satz wird umgedreht: \zitat{das ›Ich denke‹ begleitet nicht alle meine Vorstellungen, sondern alle meine Vorstellungen schliessen sich an das Ich denke an.}\footnote{Baumann, Lehren, Bd.\,2, S.\,662.}

Die zweite Zeit ist die gewöhnliche, psychologisch-astronomische. Sie ist uns gegeben an unserem Leibe und seiner Wechselwirkung mit den äußeren Dingen: Tag und Nacht, Wachen und Schlaf, die regelmäßigen Ereignisse der Natur. Die dritte ist die astronomische Zeit, die Sternzeit, die ihre Gleichförmigkeit aus der Annahme unveränderter Erdrotation gewinnt; aus ihr entsteht \zitat{jenes Idealbild der Zeit}, das gemeinhin als die Zeit schlechtweg gilt. \zitat{Dazu kam, dass man diese astronomische Zeit als eine Parallele zum Raum ansah, während beide Vorstellungen gar nicht mit einander stehen und fallen.}\footnote{Baumann, Lehren, Bd.\,2, S.\,663--667.} Kants Zeitlehre bedarf deshalb der Berichtigung: Er ging vom astronomischen Idealbild aus und wollte es als reine Anschauung gleich der des Raumes erweisen. Und die Gleichstellung von Raum und Zeit ist \zitat{durchaus unzulässig}.\footnote{Baumann, Lehren, Bd.\,2, S.\,668.}

Man kann Baumanns drei Zeiten als eine Ordnung der Parallelisierung lesen. Die dritte Zeit, das Idealbild, ist die Zeit als Linie; sie ist es, die man nach dem Raum behandelt hat. Die zweite ist die Zeit am Leib, die Zeit, die an Tag und Nacht abgelesen wird; ihr entspricht in der Physik die Zeit längs der Weltlinie eines Dinges, die Reichenbach aus der Genidentität gewinnt. Die erste ist die Zeit, die auf das Ich bezogen ist, dessen Dauer nicht Zeit ist. Je weiter man von der Linie zurückgeht, desto weniger gleicht die Zeit dem Raum.

\section{Das ewige Jetzt}

Reichenbach hat am Anfang des Zeit-Abschnitts einen Satz geschrieben, den man in einem Buch über die Physik der Raum-Zeit nicht erwartet: \zitat{›ich bin‹ ist immer gleichbedeutend mit ›ich bin jetzt‹; aber ich bin in einem ›ewigen Jetzt‹}, und fühle mich identisch bleibend.\footnote{Reichenbach, Raum-Zeit-Lehre, \S\,16.} Das Jetzt, in dem ich bin, ist immer dasselbe Jetzt, obwohl es immer ein anderes ist. Reichenbach vertagt die Erörterung dieser Ich-Zeit und kommt nicht auf sie zurück. Aber der Satz ist genau Baumanns Satz: Das Ich dauert nicht in der Zeit, sondern die Zeit wird auf das dauernde Ich bezogen. Und er steht nahe bei Brentanos Satz vom Beharren als sich erneuerndem Jetzt.

Hegel hat es in der Enzyklopädie gesagt, und Koch hat die Stelle hervorgehoben: Die Zeit ist dasselbe Prinzip wie das Ich = Ich des reinen Selbstbewusstseins, aber dasselbe in seiner Äußerlichkeit.\footnote{Hegel, Enzyklopädie \S\,258, zitiert nach Koch, Hegelsche Subjekte in Raum und Zeit.} Die Zeit ist das Ich, das sich nicht bei sich hat; das Ich ist die Zeit, die sich bei sich hat. Drei Autoren, die sich an dieser Stelle nicht aufeinander berufen, ein Physiker, ein Realist, ein Idealist, sagen dasselbe: Dass das Beharren, das die Zeit möglich macht, nicht im Raum zu suchen ist, sondern in der Dauer dessen, für den etwas jetzt ist.

\section{Das Beharrliche im Raum}

Damit ist nun Kants Umkehrung zu bedenken, die in der Widerlegung des Idealismus steht: Die Bestimmung meiner Existenz in der Zeit setzt ein Beharrliches voraus, das nicht in mir ist. Das Beharrliche ist die Substanz, und die Substanz ist im Raum. Martić hat gezeigt, wie ernst Kant das nimmt: Ohne Raum würde die Zeit selbst nicht als Größe vorgestellt werden; die Zeit wird nur am Beharrlichen zur messbaren Reihe.\footnote{Martić, Zeit und Raum, Kapitel zur Widerlegung des Idealismus und zu den Analogien.}

Die beiden Sätze, Baumanns und Kants, scheinen sich zu widersprechen: Das Beharren liegt im Ich; das Beharrliche liegt im Raum. Sie widersprechen sich nicht, wenn man zwei Fragen auseinanderhält. Kant fragt, woran die Zeit \emph{bestimmt}, das heißt gemessen und geordnet wird. Baumann fragt, woher die Zeit ihr \emph{Beharren} hat, das heißt: was es ist, das da bleibt, während die Vorstellungen wechseln. Auf die erste Frage ist die Antwort das Beharrliche im Raum, der Körper, an dem die Uhr abgelesen wird; das ist Reichenbachs Messkörper und Baumanns zweite Zeit. Auf die zweite Frage ist die Antwort das dauernde Ich. Das Beharrliche im Raum ist die äußere Gestalt der Dauer: Es ist, was die Dauer wird, wenn sie als Gegenstand angesehen wird. Brentano hat den Übergang beschrieben: Ein Ding, das besteht, beharrt durch ein sich wiederholendes Sich-Erneuern, und das ist die Weise, wie das Ich dauert, auf ein Räumliches übertragen.

Man muss hier vorsichtig sein. Es ist damit nicht gesagt, dass der Raum aus dem Ich stamme, oder das Beharrliche aus der Dauer. Es ist gesagt, dass das, was am Beharrlichen das Beharren ist, seine Weise zu sein von der Zeit hat und die Zeit die ihre vom dauernden Ich. Was der Raum ist, in dem das Beharrliche ist, bleibt dabei ungesagt. Reichenbach brauchte ihn als Messkörper, Kant als Substanz, Baumann als \zitat{Wesen sui generis}, das nicht in die Kategorientafel passt und das uns bleibt, auch wenn wir die Dinge wegdenken.\footnote{Baumann, Lehren, Bd.\,2, S.\,655; zum Raum als Anschauung, die ohne Dinge bleibt, S.\,648.} Der Raum widersteht der Ableitung an genau dieser Stelle, und die Stelle wird am Ende noch einmal aufgesucht.
